\begin{abstract}
As language models move from describing scientific results to guiding simulations and experiments, an answer can be numerically correct yet scientifically unsupported. An atomistic measurement is meaningful only together with its physical definition: which atoms were compared, in which coordinate frame, with what alignment and mathematical operation. We introduce GeoScope, a reproducible evaluation framework that tests whether language models use physically licensed measurements to justify claims about adsorbate--surface structures. Two complementary collections---55 UMA-relaxed systems and 48 OC20 DFT trajectory systems---provide a common set of measurement definitions while probing verification strength, proof selection, representation effects, and confidence under distinct evaluation protocols. In the 55-system collection, full-definition checking rejected 81 additional invalid Qwen selections that passed comparison-ID checking, including 39 correct verdicts supported by invalid proofs. The independent 48-system collection showed that these problems persist across four models: even the strongest tested configuration provided valid scientific proofs for only 567 of 960 answerable questions (59.1\%) after removing outer code fences. On the same collection, 185 of 601 highly confident Qwen3-32B selections were invalid, and no model produced a qualifying abstention on the 576 insufficient-evidence questions. These findings reveal a reliability gap that cannot be assessed by answer correctness, source identity, formatting, or token confidence alone. GeoScope establishes an inspectable foundation for testing physical evidence use in AI-assisted materials research, while leaving evidence recovery and improvements to actual autonomous discovery decisions for future validation.
\end{abstract}
